\documentclass[10pt,a4paper]{amsart}
\usepackage{amsmath,amssymb,amsthm,graphicx}
% Retained preamble for the Stacks Project chapter {CH} (shared project preamble).
% The chapter file pinned-original-source.tex only carries \input{{preamble}} on its first line,
% so no preamble text lives inside the chapter itself.  The Stacks Project uses one single shared
% preamble for every chapter: the snapshot stages it three times (sources/stacks-chow/preamble.tex,
% sources/stacks-etale-cohomology/preamble.tex and focused-seeds/AGP/sources/stacks-snapshot-20261011/
% preamble.tex) and those copies are identical in size.  The block below is copied mechanically,
% line for line, from sources/stacks-chow/preamble.tex lines 14-18 and 154-217: the Xy-pic setup and
% the native theorem declarations together with the mathematical macros.  Nothing was added, removed
% or redefined; the wrapper class and the amsmath/amssymb/amsthm/graphicx base packages are supplied
% by the assembler.
% For commutative diagrams we use Xy-pic
% Chapter: stacks-cohomology.tex -- Cohomology of Sheaves.
\usepackage[all]{xy}

% We use 2cell for 2-commutative diagrams.
\xyoption{2cell}

% Theorem environments.
%
\theoremstyle{plain}
\newtheorem{theorem}[subsection]{Theorem}
\newtheorem{proposition}[subsection]{Proposition}
\newtheorem{lemma}[subsection]{Lemma}

\theoremstyle{definition}
\newtheorem{definition}[subsection]{Definition}
\newtheorem{example}[subsection]{Example}
\newtheorem{exercise}[subsection]{Exercise}
\newtheorem{situation}[subsection]{Situation}

\theoremstyle{remark}
\newtheorem{remark}[subsection]{Remark}
\newtheorem{remarks}[subsection]{Remarks}

\numberwithin{equation}{subsection}

% Macros
%
\def\lim{\mathop{\mathrm{lim}}\nolimits}
\def\colim{\mathop{\mathrm{colim}}\nolimits}
\def\Spec{\mathop{\mathrm{Spec}}}
\def\Hom{\mathop{\mathrm{Hom}}\nolimits}
\def\Ext{\mathop{\mathrm{Ext}}\nolimits}
\def\SheafHom{\mathop{\mathcal{H}\!\mathit{om}}\nolimits}
\def\SheafExt{\mathop{\mathcal{E}\!\mathit{xt}}\nolimits}
\def\Sch{\mathit{Sch}}
\def\Mor{\mathop{\mathrm{Mor}}\nolimits}
\def\Ob{\mathop{\mathrm{Ob}}\nolimits}
\def\Sh{\mathop{\mathit{Sh}}\nolimits}
\def\NL{\mathop{N\!L}\nolimits}
\def\CH{\mathop{\mathrm{CH}}\nolimits}
\def\proetale{{pro\text{-}\acute{e}tale}}
\def\etale{{\acute{e}tale}}
\def\QCoh{\mathit{QCoh}}
\def\Ker{\mathop{\mathrm{Ker}}}
\def\Im{\mathop{\mathrm{Im}}}
\def\Coker{\mathop{\mathrm{Coker}}}
\def\Coim{\mathop{\mathrm{Coim}}}

% Boxtimes
%
\DeclareMathSymbol{\boxtimes}{\mathbin}{AMSa}{"02}

%
% Macros for moduli stacks/spaces
%
\def\QCohstack{\mathcal{QC}\!\mathit{oh}}
\def\Cohstack{\mathcal{C}\!\mathit{oh}}
\def\Spacesstack{\mathcal{S}\!\mathit{paces}}
\def\Quotfunctor{\mathrm{Quot}}
\def\Hilbfunctor{\mathrm{Hilb}}
\def\Curvesstack{\mathcal{C}\!\mathit{urves}}
\def\Polarizedstack{\mathcal{P}\!\mathit{olarized}}
\def\Complexesstack{\mathcal{C}\!\mathit{omplexes}}
% \Pic is the operator that assigns to X its picard group, usage \Pic(X)
% \Picardstack_{X/B} denotes the Picard stack of X over B
% \Picardfunctor_{X/B} denotes the Picard functor of X over B
\def\Pic{\mathop{\mathrm{Pic}}\nolimits}
\def\Picardstack{\mathcal{P}\!\mathit{ic}}
\def\Picardfunctor{\mathrm{Pic}}
\def\Deformationcategory{\mathcal{D}\!\mathit{ef}}
\usepackage{cleveref}
\crefname{originallemma}{original lemma}{original lemmas}
\Crefname{originallemma}{Original lemma}{Original lemmas}
\makeatletter
\newlabel{lemma-cech-cohomology}{{01EQ}{source}{}{sourcepointer-1}{}}
\newlabel{lemma-cech-cohomology@cref}{{[originallemma][0][]01EQ}{[1][0][]source}{}{}{}}
\makeatother
\crefname{originalequation}{original equation}{original equations}
\Crefname{originalequation}{Original equation}{Original equations}
\makeatletter
\newlabel{equation-functorial-derived}{{01F9}{source}{}{sourcepointer-2}{}}
\newlabel{equation-functorial-derived@cref}{{[originalequation][0][]01F9}{[1][0][]source}{}{}{}}
\makeatother
\crefname{originalequation}{original equation}{original equations}
\Crefname{originalequation}{Original equation}{Original equations}
\makeatletter
\newlabel{equation-functorial}{{01FA}{source}{}{sourcepointer-3}{}}
\newlabel{equation-functorial@cref}{{[originalequation][0][]01FA}{[1][0][]source}{}{}{}}
\makeatother
\crefname{originalremark}{original remark}{original remarks}
\Crefname{originalremark}{Original remark}{Original remarks}
\makeatletter
\newlabel{remark-explain-arrow}{{01FB}{source}{}{sourcepointer-4}{}}
\newlabel{remark-explain-arrow@cref}{{[originalremark][0][]01FB}{[1][0][]source}{}{}{}}
\makeatother
\crefname{originalequation}{original equation}{original equations}
\Crefname{originalequation}{Original equation}{Original equations}
\makeatletter
\newlabel{equation-choice}{{01FC}{source}{}{sourcepointer-5}{}}
\newlabel{equation-choice@cref}{{[originalequation][0][]01FC}{[1][0][]source}{}{}{}}
\makeatother
\title{Functoriality of the Cech-to-cohomology diagram along a morphism of ringed spaces}
\author{The Stacks Project Authors}
\date{stacks-project:cohomology.tex:local-snapshot-20261011; GNU FDL 1.2 or later}
\begin{document}
\maketitle
\noindent\textit{Editorial note.} Preamble: only the retained preamble of this chapter (sources/stacks-cohomology/retained-preamble.tex) is taken, because the chapter file itself carries just an input of the project's single shared preamble on its first line. It holds the Xy-pic setup, the author's native theorem declarations and the mathematical macros, copied line for line; nothing was added, removed or redefined. This card uses neither stix nor transparent, which the host TeX Live lacks. The Xy-pic matrix diagrams in the statement and in the proof rely on the xy package loaded there, which matches the two-pass local trial compilation.\par\medskip
\noindent\textit{Editorial note.} All five external references are pointers inside this chapter (one lemma, three equations, one remark) and serve to resolve the original reference commands; the arrow gamma, the four double complexes and the verification of the whole diagram are extracted verbatim inside this span and are not replaced by those pointers.\par\medskip
\noindent\textit{Editorial note.} Original numbering: in this chapter lemmas, equations and remarks share the subsection counter, and material-only assembly locates passages by label, so the printed numbers need not agree with the original ones.\par\medskip
\noindent\textit{Editorial note.} Uniform work/source\_title across every card of this source, as section 11 of NEW\_TASK\_AUTHORING.md requires (the manifest field work is taken from source\_title). The precise section of each claim is carried by claim\_key and by statement\_anchor in the registration, and the section name is repeated in the editorial notes where it matters.\par\medskip
\noindent\textit{Allowed standard prerequisites (editorial).}\begin{itemize}
\item The comparison map from the Cech complex of a covering to derived global sections. Same chapter, Lemma labelled lemma-cech-cohomology (canonical Stacks tag 01EQ, source line 1434), named in the counted statement at lines 2540 and 2542.. Supplies the horizontal arrows of the two displayed diagrams; the counted content is that the vertical arrow gamma built in the proof makes them commute.
\item Pushforward and pullback of modules along a morphism of ringed spaces, and the notion of an f-map f\textasciicircum{}*G -> F. Standard; fixed in the same chapter's section Cech cohomology and cohomology (source lines 2340 onwards) together with the displays equation-functorial-derived (tag 01F9) and equation-functorial (tag 01FA).. Fixes the objects the diagram is built from and the arrows its sides are made of.
\item Injective resolutions and double complexes computing derived global sections. Used at proof lines 2565-2591 as in the proof of Lemma lemma-cech-cohomology (tag 01EQ) and Remark remark-explain-arrow (tag 01FB, source line 2391).. Provides the four double complexes whose morphisms are compared.
\end{itemize}\medskip
\ifdefined\hypertarget\hypertarget{sourcepointer-1}{}\fi
\noindent\textit{Source reference.} Original lemma 01EQ: sources/stacks-cohomology/pinned-original-source.tex, source lines 1434--1434. Same-chapter pointer named in the counted statement (lines 2540 and 2542) and used in the proof (lines 2568 and 2619) for the horizontal arrows of the displayed diagrams, i.e. for the map from the Cech complex to derived global sections. It fixes the arrows that the counted diagram compares; no part of the comparison is taken from it.\par\medskip
\ifdefined\hypertarget\hypertarget{sourcepointer-2}{}\fi
\noindent\textit{Source reference.} Original equation 01F9: sources/stacks-cohomology/pinned-original-source.tex, source lines 2370--2370. Same-chapter pointer to the display naming the right vertical arrow of the first diagram in the counted statement (line 2541). Navigation to a display only.\par\medskip
\ifdefined\hypertarget\hypertarget{sourcepointer-3}{}\fi
\noindent\textit{Source reference.} Original equation 01FA: sources/stacks-cohomology/pinned-original-source.tex, source lines 2386--2386. Same-chapter pointer to the display giving the right vertical arrow of the cohomology diagram in the second half of the counted statement (line 2552). Navigation to a display only.\par\medskip
\ifdefined\hypertarget\hypertarget{sourcepointer-4}{}\fi
\noindent\textit{Source reference.} Original remark 01FB: sources/stacks-cohomology/pinned-original-source.tex, source lines 2391--2391. Same-chapter pointer used at proof lines 2578-2582, where the auxiliary injective resolution f\_*I -> (J')\textasciicircum{}bullet and the morphism beta are chosen exactly as explained in that remark. It points at a construction recipe; the counted proof still has to keep the four double complexes coherent.\par\medskip
\ifdefined\hypertarget\hypertarget{sourcepointer-5}{}\fi
\noindent\textit{Source reference.} Original equation 01FC: sources/stacks-cohomology/pinned-original-source.tex, source lines 2405--2405. Same-chapter pointer used at proof lines 2582 and 2621: commutativity of the pentagon in the displayed diagram is derived from commutativity of this chosen square. Navigation to the displayed commutativity that the construction must satisfy.\par\medskip
\begin{lemma}
\label{lemma-functoriality-cech}
Let $f : X \to Y$ be a morphism of ringed spaces.
Let $\varphi : f^*\mathcal{G} \to \mathcal{F}$ be an $f$-map
from an $\mathcal{O}_Y$-module $\mathcal{G}$ to an
$\mathcal{O}_X$-module $\mathcal{F}$.
Let $\mathcal{U} : X = \bigcup_{i \in I} U_i$ and
$\mathcal{V} : Y = \bigcup_{j \in J} V_j$ be open coverings.
Assume that $\mathcal{U}$ is a refinement of
$f^{-1}\mathcal{V} : X = \bigcup_{j \in J} f^{-1}(V_j)$.
In this case there exists a commutative diagram
$$
\xymatrix{
\check{\mathcal{C}}^\bullet(\mathcal{U}, \mathcal{F}) \ar[r] &
R\Gamma(X, \mathcal{F}) \\
\check{\mathcal{C}}^\bullet(\mathcal{V}, \mathcal{G}) \ar[r]
\ar[u]^\gamma &
R\Gamma(Y, \mathcal{G}) \ar[u]
}
$$
in $D^{+}(\mathcal{O}_X(X))$ with horizontal arrows given by
Lemma \ref{lemma-cech-cohomology} and right vertical arrow by
(\ref{equation-functorial-derived}).
In particular we get commutative diagrams of cohomology groups
$$
\xymatrix{
\check{H}^p(\mathcal{U}, \mathcal{F}) \ar[r] &
H^p(X, \mathcal{F}) \\
\check{H}^p(\mathcal{V}, \mathcal{G}) \ar[r]
\ar[u]^\gamma &
H^p(Y, \mathcal{G}) \ar[u]
}
$$
where the right vertical arrow is (\ref{equation-functorial})
\end{lemma}

\begin{proof}
We first define the left vertical arrow. Namely, choose a map
$c : I \to J$ such that $U_i \subset f^{-1}(V_{c(i)})$ for all
$i \in I$. In degree $p$ we define the map by the rule
$$
\gamma(s)_{i_0 \ldots i_p} = \varphi(s)_{c(i_0) \ldots c(i_p)}
$$
This makes sense because $\varphi$ does indeed induce maps
$\mathcal{G}(V_{c(i_0) \ldots c(i_p)}) \to \mathcal{F}(U_{i_0 \ldots i_p})$
by assumption. It is also clear that this defines a morphism of complexes.
Choose injective resolutions
$\mathcal{F} \to \mathcal{I}^\bullet$ on $X$ and
$\mathcal{G} \to J^\bullet$ on $Y$. According to
the proof of Lemma \ref{lemma-cech-cohomology} we introduce the double
complexes $A^{\bullet, \bullet}$ and $B^{\bullet, \bullet}$
with terms
$$
B^{p, q} = \check{\mathcal{C}}^p(\mathcal{V}, \mathcal{J}^q)
\quad
\text{and}
\quad
A^{p, q} = \check{\mathcal{C}}^p(\mathcal{U}, \mathcal{I}^q).
$$
As in Remark \ref{remark-explain-arrow} above we also choose an
injective resolution
$f_*\mathcal{I} \to (\mathcal{J}')^\bullet$ on $Y$ and a morphism
of complexes $\beta : \mathcal{J} \to (\mathcal{J}')^\bullet$
making (\ref{equation-choice}) commutes. We introduce some more
double complexes, namely $(B')^{\bullet, \bullet}$ and
$(B'')^{\bullet, \bullet}$ with
$$
(B')^{p, q} = \check{\mathcal{C}}^p(\mathcal{V}, (\mathcal{J}')^q)
\quad
\text{and}
\quad
(B'')^{p, q} = \check{\mathcal{C}}^p(\mathcal{V}, f_*\mathcal{I}^q).
$$
Note that there is an $f$-map of complexes from
$f_*\mathcal{I}^\bullet$ to $\mathcal{I}^\bullet$. Hence
it is clear that the same rule as above defines a morphism
of double complexes
$$
\gamma : (B'')^{\bullet, \bullet} \longrightarrow A^{\bullet, \bullet}.
$$
Consider the diagram of complexes
$$
\xymatrix{
\check{\mathcal{C}}^\bullet(\mathcal{U}, \mathcal{F})
\ar[r] &
\text{Tot}(A^{\bullet, \bullet}) & & &
\Gamma(X, \mathcal{I}^\bullet) \ar[lll]^{qis}
\ar@{=}[ddl]\\
\check{\mathcal{C}}^\bullet(\mathcal{V}, \mathcal{G})
\ar[r] \ar[u]^\gamma &
\text{Tot}(B^{\bullet, \bullet}) \ar[r]^\beta &
\text{Tot}((B')^{\bullet, \bullet}) &
\text{Tot}((B'')^{\bullet, \bullet}) \ar[l] \ar[llu]_{s\gamma} \\
& \Gamma(Y, \mathcal{J}^\bullet) \ar[u]^{qis} \ar[r]^\beta &
\Gamma(Y, (\mathcal{J}')^\bullet) \ar[u] &
\Gamma(Y, f_*\mathcal{I}^\bullet) \ar[u] \ar[l]_{qis}
}
$$
The two horizontal arrows with targets $\text{Tot}(A^{\bullet, \bullet})$ and
$\text{Tot}(B^{\bullet, \bullet})$
are the ones explained in Lemma \ref{lemma-cech-cohomology}.
The left upper shape (a pentagon) is commutative simply
because (\ref{equation-choice}) is commutative.
The two lower squares are trivially commutative.
It is also immediate from the definitions that the
right upper shape (a square) is commutative.
The result of the lemma now follows from the definitions
and the fact that going around the diagram on the outer sides
from $\check{\mathcal{C}}^\bullet(\mathcal{V}, \mathcal{G})$
to $\Gamma(X, \mathcal{I}^\bullet)$ either on top or on bottom
is the same (where you have to invert any quasi-isomorphisms along the way).
\end{proof}

\end{document}
